\chapter*{Summary}
\addcontentsline{toc}{chapter}{Summary}

MASFE (Multi-Algorithm Scheduling and Fusion Engine) is an Orbital Edge AI system designed for the NASA Space-to-Soil Challenge 2026. It directly addresses NASA's stated objective: to orchestrate existing land observation algorithms into an efficient, responsive onboard intelligence layer for regenerative agriculture SmallSat missions.

A reinforcement-learning-inspired Markov Decision Process (MDP) scheduler runs on a radiation-tolerant LEON4FT CPU, dynamically allocating compute across two NASA algorithms---MODIS MOD13A1 (Enhanced Vegetation Index, DOI: 10.5067/MODIS/MOD13A1.061) and MOD11A1 (Land Surface Temperature, DOI: 10.5067/MODIS/MOD11A1.061)---fused onboard via an FPGA preprocessing stage into a per-field Crop Stress Composite (CSC). The satellite makes \emph{real-time, on-orbit decisions} about where to sense, which algorithms to execute, and what to downlink.

\medskip
\noindent Key results from simulation, compared to a naive raw-dump baseline:

\begin{itemize}
  \item \textbf{97\%} reduction in downlink data volume (14{,}400\,MB $\rightarrow$ 423\,MB per 120-pass season)
  \item \textbf{100\%} pathogen detection retention
  \item \textbf{$<$90\,min} alert latency versus 1--3 days for ground processing
  \item \textbf{9.6\,W} peak power --- within the 6U CubeSat 14\,W average budget
\end{itemize}

\medskip
\noindent Deployed as a hosted payload on a Loft Orbital YAM platform and integrated with the OpenET API, MASFE delivers sub-90-minute crop pathogen risk alerts to regenerative farmers. Early detection enables targeted fungicide application against the global 8-billion-pounds-per-year pesticide burden documented by Dr.\ Katie Gold (Cornell University, NASA Webinar~1, 2026), addressing a \$220\,billion annual crop loss problem.

The MDP orchestration engine is domain-agnostic---agriculture is the demonstration; the scheduler is the contribution and licensable IP. The system targets the NASA AIST TRL~4 hardware demonstration pathway as its immediate next step.
